
% !TEX root = ../../../Esperimenti/.tex/MEF.tex
% LTeX: language=it

\begin{tikzpicture}
    %\begingroup Parametri
        % \renewcommand\labelSize{10}
        % \renewcommand\tickSize{9}
        % \renewcommand\legendSize{9}
    
        \newcommand\width{10cm}
        \pgfmathsetlengthmacro\height{0.618*\width}
    %\endgroup

    %\begingroup Funzione lognormale
        \pgfmathsetmacro\mi{1}
        \newcommand\plotLognormal[1]{
            \foreach \s in {2.5,2,1.5,1,0.5,0.3,0.2,0.13,0.1}{%
                \addplot+ [thick,no marks] gnuplot [raw gnuplot] {%
                    set samples 5000;
                    m=1; s=\s;
                    f(x,m,s)=(1/(x*s*sqrt(2*pi)))*exp(-(log(x)-m)**2/(2*s**2));
                    #1
                };

                \addlegendentryexpanded{\(\sigma=\s\)}
            }

            \addplot [draw=none,forget plot] gnuplot [raw gnuplot] {%
                f(x,m,s)=(1/(x*s*sqrt(2*pi)))*exp(-(log(x)-m)**2/(2*s**2));
                m=1; s=0.1; plot[x=exp(\mi):exp(\mi)] (f(x,m,s));
            } coordinate (media); 
            \draw [thick,dashed,black] (axis cs:{exp(\mi)},0) -- (media);
        }
    %\endgroup

    %\begingroup Impostazioni
        \pgfplotsset{
            /pgfplots/group/goupStyle/.style={
                group name=Lognormale,
                vertical sep=2em,
                % x descriptions at=edge bottom,
                y descriptions at=edge left,
                group size=1 by 2,
            },
            plotStyle/.style={
                width=\width,
                height=\height,
                % x post scale=0.7,
                % y post scale=0.8,
                scale only axis,
                axis lines=left,
                ymin=0,ymax=1.5,
                xlabel={\(x\)},ylabel={\(f_X(x)\)},
                xlabel style={font=\dimTesto{\labelSize}\color{white!15!black}},
                ylabel style={font=\dimTesto{\labelSize}\color{white!15!black}},
                xticklabel style={font=\dimTesto{\tickSize}},
                yticklabel style={font=\dimTesto{\tickSize}},
                legend style={
                    legend pos=north east,
                    legend cell align=left,
                    font=\dimTesto{\legendSize},
                    draw=none,
                    fill=white,
                    fill opacity=.5,
                    text opacity=1,
                    % nodes={
                    %     scale=0.8,
                    %     transform shape
                    % }
                },
                grid=major,
                grid style={
                    dashed,
                    gray!20
                },
                colormap/viridis,
                cycle list={[samples of colormap={9} of viridis]},
            }
        }
    %\endgroup

    \begin{groupplot}[group style={goupStyle},plotStyle]
        %\begingroup
            \gdef\xmin{1e-10}
            \gdef\xmax{6}
            \nextgroupplot[
                xmin=\xmin,xmax=\xmax,
                xmode=linear,xlabel={},
            ]\plotLognormal{plot [x=\xmin:\xmax] f(x,m,s);}
            %\legend{}
        %\endgroup

        %\begingroup
            \gdef\xmin{1e-5}\gdef\xmax{10}
            \nextgroupplot[
                xmin=1e-5,xmax=\xmax,
                xmode=log,legend style={
                    legend pos=north west,
                },
            ]\plotLognormal{
                set parametric;
                plot [x=log(\xmin):log(\xmax)] exp(x), f(exp(x),m,s);
            }
            %\legend{}
        %\endgroup
    \end{groupplot}

    % \redefineTikZbounds{Lognormale c1r1}{Lognormale c1r2}
\end{tikzpicture}

%\begingroup Vecchia versione
    % \begin{tikzpicture}%
    %     % La lognormale è «f(x)=(1/(x*sigma*sqrt(2*pi)))*exp(-(ln(x)-mu)^2/(2*sigma^2))»
    
    %     \newcommand\muPrm{0}    % Media della normale sottostante
    %     \newcommand\sigmaPrm{1} % Deviazione standard della normale sottostante
    
    %     % Definizione della funzione matematica pgf
    %     \pgfmathdeclarefunction{lognormalpdf}{1}{%
    %         \pgfmathparse{1/(#1*\sigmaPrm*sqrt(2*pi))*exp(-((ln(#1)-\muPrm)^2)/(2*\sigmaPrm*\sigmaPrm))}%
    %     }
    
    %     % \renewcommand\subCaptionSize{10}
    %     % \renewcommand\labelSize{10}
    %     % \renewcommand\tickSize{9}
    
    %     \pgfmathsetmacro\hsep{1cm}
    %     % \pgfmathsetmacro\plotWidth{(\linewidth-\hsep)/2}
    %     \pgfmathsetmacro\plotWidth{4.8cm}
    %     \pgfmathsetmacro\plotHeight{.618*\plotWidth}
    
    %     \begin{groupplot}[
    %         group style={
    %             group name=lgnrml,
    %             group size=2 by 1,
    %             horizontal sep=\hsep,
    %             xlabels at=edge bottom,
    %             ylabels at=edge left,
    %             x descriptions at=edge bottom,
    %             y descriptions at=edge left
    %         },
    %         width=\plotWidth,
    %         height=\plotHeight,
    %         scale only axis,
    %         axis x line*=bottom,
    %         axis y line*=left,
    %         xlabel={$x$},
    %         ylabel={$f_X(x)$},
    %         xlabel style={font=\dimTesto{\labelSize}\color{white!15!black}},
    %         ylabel style={font=\dimTesto{\labelSize}\color{white!15!black}},
    %         xticklabel style={font=\dimTesto{\tickSize}},
    %         yticklabel style={font=\dimTesto{\tickSize}},
    %         grid=major,
    %         grid style={dashed},
    %         every axis plot/.append style={
    %             % ultra thick,
    %             thick,
    %             smooth
    %         },
    %         ymin=0,
    %     ]
    %         % Scala lineare
    %         \nextgroupplot[
    %             xmin=0,
    %             xmax=10, %5
    %             minor tick num=1,
    %         ]
    %         \addplot[
    %             % color=blue,
    %             domain=0.01:20,
    %             samples=400
    %         ] {lognormalpdf(x)};
    
    %         % Scala semilogaritmica
    %         \nextgroupplot[
    %             xmode=log,
    %             xmin=0.01,
    %             xmax=20,
    %             % log ticks with fixed point, % Mostra i numeri (0.1,1,10) invece di (10^-1,10^0,10^1)
    %         ]
    %         % Nota: il dominio è molto più ampio qui per mostrare le code
    %         \addplot[
    %             % color=red,
    %             domain=0.005:25,
    %             samples=500
    %         ] {lognormalpdf(x)};
    %     \end{groupplot}
    
    %     \tikzset{SubCaption/.style={
    %         text width=\plotWidth,
    %         anchor=north,
    %         align=center,
    %         yshift=-2.5em
    %     }}
    %     \node[SubCaption] at (lgnrml c1r1.south)
    %         {\dimTesto{\subCaptionSize}\sottoDidascalia{Scala lineare.\label{figLognormaleScalaLineare}}};
    %     \node[SubCaption] at (lgnrml c2r1.south)
    %         {\dimTesto{\subCaptionSize}\sottoDidascalia{Scala semilogaritmica.\label{figLognormaleScalaSemiLogaritmica}}};
    
    %     \redefineTikZbounds{lgnrml c1r1}{lgnrml c2r1}
    % \end{tikzpicture}%
%\endgroup
